\chapter{Literature Review}
\label{ch:literature}

\section{Overview}
\label{sec:lit-overview}

The software landscape surrounding professional sport is broad but fragmented. Existing
products tend to specialise narrowly: some excel at delivering live scores and statistics to a
mass audience, others at capturing biometric performance data for elite teams, and others
again at the logistical chores of amateur club administration. No single category addresses
the full operational breadth of a multi-sport club, and almost none combine consumer-grade
real-time data with the internal, role-aware workflows that staff require. This chapter
surveys the principal categories of relevant systems, examines their strengths and the gaps
they leave open, and concludes by positioning MSCMS (Sportify) against representative
platforms.

The review is organised into two complementary parts. The first part is a functional survey
of existing commercial systems, grouped into three families: large-audience sports-data
platforms, elite performance and analysis suites, and generic club- and team-administration
tools. For each family the survey identifies what the products do well, what they omit, and
why those omissions matter to a multi-sport club. The second part steps back from individual
products to review the technical literature on the architectural principles that MSCMS
relies upon---microservices decomposition, role-based access control, event-driven
integration and real-time web communication---so that the design decisions taken later in
this report can be traced to established practice rather than presented as ad hoc choices.
The chapter then synthesises both parts into an explicit statement of the research gap and
the positioning of the proposed system, supported by a multi-dimensional comparison.

\section{Sports-Data and Live-Score Platforms}
\label{sec:lit-data}

Consumer-facing sports-data platforms such as SofaScore and ESPN have defined modern
expectations for how sporting information is presented. Their strengths are considerable.
They aggregate fixtures, results, league standings and rich in-match event timelines across a
vast range of competitions, and they deliver this information in real time, updating scores
and statistics live as matches unfold. They cover dozens of sports and hundreds of
competitions, normalising heterogeneous data into a consistent presentation, and they have
invested heavily in responsive, mobile-first interfaces that render comfortably on any device.
The underlying data is sourced from professional feeds, and a comparable feed---exemplified
by the football-data.org API~\cite{footballdata}---exposes competitions, fixtures, results and
standings through a clean RESTful interface~\cite{restapi}, allowing third parties to build on
authoritative match data rather than re-keying it by hand. These platforms have
demonstrated, at very large scale, that supporters value immediacy, breadth of coverage and
a polished presentation of live events, and they have established the architectural pattern of
ingesting authoritative third-party data and re-serving it to clients through a fast read path.

Their weaknesses, from the standpoint of a club operator, are equally clear. They are
fundamentally read-only, outward-facing products: they tell a global audience what is
happening on the pitch, but they hold no notion of an individual club's internal operations.
They model neither rosters and contracts nor training, medical or scouting workflows; they
have no concept of staff roles or fine-grained access control beyond a generic consumer
account; and they offer no mechanism for a club to record its own line-ups, run its own
competitions, or coordinate its own staff. A club cannot enter the result of a youth-team
friendly into ESPN, nor restrict a salary figure to its commercial department within
SofaScore, because these systems were never designed to be systems of record for a single
organisation. In short, they are excellent sources of match data and exemplary models of
real-time presentation, but they are consumption platforms rather than management
platforms. MSCMS deliberately adopts their real-time presentation philosophy and ingests
the same class of authoritative data~\cite{footballdata}---storing it internally so that the
backend, not the external feed, is the single source of truth---while extending far beyond
consumption into the internal management workflows these platforms omit entirely.

\section{Elite Performance and Analysis Suites}
\label{sec:lit-performance}

A second family of products targets the performance department of professional teams.
Video-analysis and performance platforms such as Hudl provide tools for tagging match
footage, sharing clips with players and building analytical reports, and have become a
fixture of professional coaching workflows from elite academies to senior squads.
Athlete-monitoring vendors such as Catapult specialise in capturing biometric and load data
from wearable devices---accelerometers, gyroscopes and positioning sensors---supporting
fitness staff in managing training load and reducing injury risk through quantified
workload metrics. Enterprise offerings, of which SAP Sports One is representative, integrate
squad management, training planning and some medical record-keeping into a single
commercial suite aimed at well-resourced clubs, and demonstrate that several domains can be
brought under one commercial roof when sufficient engineering investment is made.

These suites are powerful within their chosen niche. Hudl's video tooling and Catapult's
load monitoring are genuinely best-in-class, and the enterprise suites prove that domain
integration is achievable. However, each leaves significant gaps relative to the needs of a
multi-sport club. Their multi-sport modelling is typically limited: they assume a single
discipline, or treat additional sports as loosely coupled add-ons rather than first-class
entities with their own positions, formats and rules. Their role models are coarse, oriented
towards a handful of coaching and performance roles rather than the full spectrum of
leadership, medical, scouting, commercial and supporter roles that a club actually contains;
they have no place for a sponsor negotiating an offer or a supporter following the club. They
are predominantly closed, proprietary and cloud-hosted by the vendor, which means a club
cannot self-host them, inspect their internals, audit their handling of sensitive medical data,
or adapt them to local needs. They rarely expose a coherent fan-facing surface or an
embedded predictive-analytics capability that non-technical staff can consume directly, and
where they do offer analytics, it is confined to the performance domain and does not span
medical, commercial and competition data within a single store. The result is depth without
breadth: a club that adopts such a suite still finds its competition records, its commercial
pipeline and its supporter engagement living elsewhere.

\section{Generic Club and Team Administration Tools}
\label{sec:lit-admin}

A third family addresses the administrative and logistical needs of clubs and teams,
particularly at the amateur and semi-professional levels. Tools such as TeamSnap focus on
scheduling, attendance, team communication and payment collection, and they are widely
adopted because they are inexpensive, easy to use and effective at the chores they target.
Their strength is accessibility: a volunteer organiser can have a team running on such a tool
within minutes, with no technical expertise and no infrastructure to maintain. For their
intended audience this low barrier to entry is precisely the point.

Their limitations, however, make them unsuitable as the operational backbone of a
professional multi-sport club. They model a single team rather than a multi-discipline
governing body; their notion of roles seldom extends beyond organiser, player and parent;
and they provide no structured medical workflow, no scouting or transfer pipeline, no
sponsorship management and no embedded analytics or machine learning. They are not
designed to ingest authoritative competition data, to run a live-match engine that drives a
fixture through its phases as it is played, or to enforce the layered, fine-grained access
control~\cite{rbac} that a professional organisation requires to protect sensitive medical and
financial information. As with the performance suites, they are closed and vendor-hosted,
offering no path to self-hosting or architectural extension. They are, in effect, the inverse of
the enterprise suites: broad in their accessibility but shallow in every domain a professional
club depends upon.

\section{Technical Foundations in the Literature}
\label{sec:lit-technical}

Beyond surveying competing products, it is necessary to ground the architecture of MSCMS
in the established software-engineering literature, since the value of the system lies as much
in how it is built as in what it does. Four bodies of work are particularly relevant.

\subsection{Microservices Architecture}
\label{subsec:lit-micro}

The microservices style, articulated influentially by Lewis and Fowler~\cite{microservices},
advocates building an application as a suite of small, independently deployable services, each
running in its own process, organised around business capabilities, and communicating
through lightweight network mechanisms. The literature identifies several benefits directly
relevant to a multi-sport club: independent deployability allows each domain to evolve and
release on its own cadence; fault isolation contains a failure within one service rather than
propagating it across the whole system; and the freedom to choose technology and scaling
strategy per service allows a bursty real-time workload to be treated differently from a
steady transactional one. The same literature is candid about the costs---distributed-system
complexity, the need for service discovery and centralised configuration, and the discipline
required to keep service boundaries aligned with data ownership. MSCMS adopts the
database-per-service pattern that the literature recommends as the means of preserving loose
coupling, and it employs a discovery server, a configuration server and an API gateway from
the Spring Cloud~\cite{springcloud} ecosystem to manage exactly the operational concerns
the literature warns of, building on the Spring Boot~\cite{springboot} platform.

\subsection{Role-Based Access Control}
\label{subsec:lit-rbac}

Role-based access control, formalised by Ferraiolo, Kuhn and Chandramouli~\cite{rbac},
mediates a user's permissions through the roles they hold rather than through permissions
assigned to individuals directly. This indirection is well suited to an organisation with many
job functions and a high turnover of personnel: permissions are defined once per role, and
onboarding or off-boarding a member becomes a matter of assigning or revoking a role
rather than reconfiguring dozens of individual rights. The literature emphasises that access
control is only as strong as its weakest enforcement point; a control that exists solely in a
user interface is no control at all, because the interface can be bypassed. MSCMS therefore
implements RBAC at three layers---the identity provider, the API gateway and the front-end
route guards---over a catalogue of fourteen roles, so that a request for an unauthorised
resource is denied regardless of the path by which it arrives. Identity itself is delegated to a
mature, standards-based provider~\cite{keycloak} that issues signed JSON Web
Tokens~\cite{jwt}, following the well-understood OAuth~2.0 and OpenID~Connect patterns
rather than a bespoke scheme.

\subsection{Event-Driven Architecture and Messaging}
\label{subsec:lit-eda}

Event-driven architecture decouples the producer of an event from its consumers: a service
publishes a fact about something that has happened, and any number of other services may
react to it without the publisher knowing or caring who they are. In a distributed system this
decoupling is the natural complement to microservices, because it lets domains integrate
without the tight, synchronous coupling that would otherwise re-create a monolith over the
network. The literature on durable, partitioned log-based brokers, of which Apache
Kafka~\cite{kafka} is the canonical example, describes how events can be published reliably
and consumed by independent subscribers at their own pace. MSCMS uses this pattern
directly: a goal recorded by the live-match engine, an injury logged by the medical service or
a transfer concluded by the commercial workflow is published as an event, and the
notification service consumes those events to raise alerts---so that the alerting logic lives
entirely outside the domains that generate the events. The outbox pattern is employed to
ensure that an event is published if and only if the corresponding database change commits,
avoiding the classic inconsistency between a service's state and the events it has emitted.

\subsection{Real-Time Web Communication}
\label{subsec:lit-rt}

Finally, real-time features such as live match updates and team chat require a communication
mechanism richer than the request--response model of classical HTTP. The WebSocket
protocol~\cite{websocket} provides a full-duplex, persistent connection over a single TCP
socket, allowing a server to push data to a client the instant it becomes available rather
than waiting to be polled. MSCMS uses a WebSocket channel for its real-time team chat,
with a polling fallback for environments in which a persistent connection cannot be
established, reflecting the established engineering practice of degrading gracefully rather
than failing outright when the optimal transport is unavailable.

\section{Architectural Perspective on Existing Products}
\label{sec:lit-arch}

Returning to the surveyed products with these foundations in mind, the existing commercial
systems share an architectural characteristic worth noting: they are, almost without
exception, closed monolithic services hosted exclusively by their vendors. This delivery model
has concrete consequences for clubs. It prevents self-hosting and on-premises data residency,
which matters acutely for sensitive medical and financial records subject to data-protection
obligations; it couples every functional area to a single vendor release cadence, so that a
club cannot evolve one capability without waiting on the vendor's roadmap; and it offers no
opportunity to extend the platform or integrate it deeply with a club's own systems. The
software-engineering literature on microservices~\cite{microservices} argues persuasively that
decomposing a large system into independently deployable services improves fault isolation,
scalability and the ability to evolve individual domains without destabilising the whole, and
the literature on RBAC~\cite{rbac} and event-driven integration~\cite{kafka} supplies the
remaining pillars. MSCMS adopts this architectural stance explicitly: by building on a
microservices backend with a self-hosted identity provider~\cite{keycloak} and a
self-hostable, containerised~\cite{docker} deployment, it offers a degree of openness,
extensibility and data ownership that the surveyed commercial products do not.

\section{Comparative Positioning}
\label{sec:lit-comparison}

Table~\ref{tab:platform-comparison} positions MSCMS (Sportify) against representative
platforms drawn from the three families surveyed above: SofaScore and ESPN as
large-audience data platforms, Hudl and Catapult as performance suites, SAP Sports One as
an enterprise suite, and TeamSnap as a generic administration tool. The comparison is
deliberately multi-dimensional, spanning functional capabilities (live-match engine, data
ingestion, medical and commercial workflows, messaging and analytics), architectural
properties (microservices, self-hosted identity, self-hostable deployment) and audience
reach (fan-facing surface). For brevity the consumer data platforms are grouped into a single
column, since their capabilities with respect to internal club management are uniformly
absent, as are the enterprise suite and the administration tool grouped where their profiles
coincide. The comparison highlights that each existing product is strong within its niche but
leaves substantial gaps elsewhere, whereas MSCMS is designed specifically to span all of the
surveyed dimensions in one coherent system.

\begin{table}[htbp]
  \centering
  \caption{Multi-dimensional comparison of MSCMS (Sportify) against representative sports
    platforms across functional, architectural and audience dimensions.
    \emph{Yes} denotes first-class support, \emph{Part.} partial or domain-limited support,
    and \emph{No} the absence of the capability.}
  \label{tab:platform-comparison}
  \small
  \begin{tabular}{@{}lccccc@{}}
    \toprule
    \textbf{Capability} & \textbf{SofaScore/} & \textbf{Hudl} & \textbf{Catapult} &
      \textbf{SAP/} & \textbf{MSCMS} \\
    & \textbf{ESPN} & & & \textbf{TeamSnap} & \textbf{(Sportify)} \\
    \midrule
    \multicolumn{6}{@{}l}{\emph{Functional capabilities}}\\
    Real-time live-match engine      & Part. & No   & No   & No    & Yes \\
    Real competition-data ingestion  & Yes   & No   & No   & Part. & Yes \\
    Competition standings \& brackets& Yes   & No   & No   & Part. & Yes \\
    Multi-sport modelling            & Part. & Part.& Part.& Part. & Yes \\
    Squad, contracts \& transfers    & No    & Part.& No   & Part. & Yes \\
    Structured medical workflow      & No    & No   & Part.& Part. & Yes \\
    Scouting pipeline                & No    & Part.& No   & Part. & Yes \\
    Sponsorship workflow             & No    & No   & No   & Part. & Yes \\
    Real-time team messaging         & No    & Part.& No   & Yes   & Yes \\
    Event-driven notifications       & Part. & Part.& Part.& Part. & Yes \\
    Embedded ML predictions          & No    & No   & No   & Part. & Yes \\
    \midrule
    \multicolumn{6}{@{}l}{\emph{Architectural properties}}\\
    Fine-grained role-based access   & No    & Part.& Part.& Part. & Yes \\
    Microservices architecture       & No    & No   & No   & Part. & Yes \\
    Self-hosted identity provider    & No    & No   & No   & No    & Yes \\
    Self-hostable deployment         & No    & No   & No   & No    & Yes \\
    \midrule
    \multicolumn{6}{@{}l}{\emph{Audience reach}}\\
    Internal staff workflows         & No    & Yes  & Yes  & Yes   & Yes \\
    Fan-facing surface               & Yes   & No   & No   & No    & Yes \\
    \bottomrule
  \end{tabular}
\end{table}

The table makes the pattern visible at a glance. The consumer data platforms occupy the top
of the audience-reach dimension---they are unmatched at fan-facing presentation and data
ingestion---but are essentially blank across the internal-workflow and architectural rows.
The performance and enterprise suites are the mirror image: strong on internal staff
workflows but weak or absent on fan reach, real competition data and self-hostable, open
architecture. The generic administration tool is accessible but shallow across nearly every
functional dimension. MSCMS is the only column that registers first-class support across all
three dimensions, which is precisely the claim the system sets out to substantiate: not that it
beats any specialist product within that specialist's own niche, but that it is the only
platform that unifies the niches into a single, coherent, self-hostable whole.

\section{Research Gap and Positioning}
\label{sec:lit-gap}

The survey reveals a consistent pattern that constitutes the research gap this project
addresses. Large-audience data platforms have mastered real-time presentation and
authoritative match data but hold no notion of a club's internal operations. Performance and
enterprise suites are deep within a single domain but narrow in role coverage, weak in
multi-sport modelling, and closed to self-hosting and extension. Generic administration tools
are accessible but shallow, modelling a single team without medical, scouting, sponsorship or
analytical capability. Across all three families, fine-grained role-based access
control~\cite{rbac}, an open microservices architecture~\cite{microservices}, a self-hosted
identity provider~\cite{keycloak}, a structured medical workflow and an embedded predictive
capability are conspicuously rare or entirely absent, and---most tellingly---no single product
unifies real-time, authoritative match data~\cite{footballdata} with the full internal
management of a multi-discipline club.

The gap, stated precisely, is the absence of a system that is simultaneously (i)
\emph{comprehensive}, spanning identity, squad, competitions, live matches, medical,
commercial, communication and analytics; (ii) \emph{real-time and event-driven}, driving
live matches and propagating cross-domain events~\cite{kafka} rather than recording facts
after the event; (iii) \emph{role-aware} through layered, enforced access control across
fourteen distinct professional roles; (iv) \emph{open and self-hostable}, built on a
microservices backend~\cite{springboot,springcloud} that a club can run and extend itself;
and (v) \emph{dual-audience}, serving internal staff and external supporters from the same
authoritative data store. Each surveyed product satisfies a subset of these properties; none
satisfies all of them together, and it is the act of satisfying them together---reconciling
consistent identity, consistent access control and consistent data semantics across
heterogeneous domains---that constitutes the substantive contribution.

This gap is precisely the opportunity that MSCMS (Sportify) addresses. By combining the
real-time, data-rich presentation philosophy of the large-audience platforms with the
internal-management depth of the performance and administration tools, and by delivering
the whole as an open, self-hostable microservices system with layered access control over a
self-hosted identity provider, MSCMS offers a class of platform that the existing market does
not provide. The chapters that follow describe the methodology by which this system was
designed and built, grounding each architectural decision in the principles reviewed above,
and the results that demonstrate it in operation.
